% ==============================================================================
% MÓDULO: ANEXOS OFICIALES (ENFOQUE CUALITATIVO FENOMENOLÓGICO)
% Archivo: subcapitulos/sec_07_anexos.tex
% ==============================================================================
\appendix
\section*{ANEXOS}
\label{sec:anexos}
\addcontentsline{toc}{section}{ANEXOS}

% ------------------------------------------------------------------------------
% ANEXO 1: GUÍA DE ENTREVISTA A PROFUNDIDAD FENOMENOLÓGICA
% ------------------------------------------------------------------------------
\phantomsection
\subsection*{Anexo 1: Guía de Entrevista a Profundidad Fenomenológica Semiestructurada Bilingüe}
\label{anx:1_guia}
\addcontentsline{toc}{subsection}{Anexo 1: Guía de Entrevista a Profundidad Fenomenológica}

\noindent\textbf{Título de la investigación:} Percepción de la calidad de atención en usuarios del Centro de Salud Belén, Ayacucho 2026.\\
\textbf{Investigadora:} Gresly Lucero Pariona Palomino \quad | \quad \textbf{Asesor:} Dr. Manglio Aguirre Andrade.\\
\textbf{Propósito cualitativo:} Explorar, describir y comprender la esencia de la experiencia vivida (\textit{vivencias}) por el usuario respecto al acceso y la calidad del trato humano y cuidado en el Centro de Salud Belén.

\vspace{0.3cm}
\noindent\textbf{I. Ficha de Caracterización del Informante Clave (Muestra de Casos Tipo):}
\begin{itemize}
    \item Código del participante: \line(1,0){60} \quad | \quad Fecha: \line(1,0){50} \quad | \quad Hora inicio/fin: \line(1,0){50}
    \item Edad: \line(1,0){25} años \quad | \quad Sexo: [ \phantom{X} ] Femenino \quad [ \phantom{X} ] Masculino
    \item Procedencia (Barrio / Sector): \line(1,0){120} \quad | \quad Tiempo viviendo en Huamanga: \line(1,0){30}
    \item Idioma materno: [ \phantom{X} ] Quechua chanka \quad [ \phantom{X} ] Castellano \quad [ \phantom{X} ] Bilingüe
    \item Nivel de instrucción: [ \phantom{X} ] Sin instrucción \quad [ \phantom{X} ] Primaria \quad [ \phantom{X} ] Secundaria \quad [ \phantom{X} ] Superior
    \item Ocupación principal: \line(1,0){100} \quad | \quad Tipo de seguro: [ \phantom{X} ] SIS subsidiado \quad [ \phantom{X} ] Otro
    \item Servicio asistencial donde recibió atención hoy: [ \phantom{X} ] Medicina \quad [ \phantom{X} ] CRED \quad [ \phantom{X} ] Inmunizaciones \quad [ \phantom{X} ] Obstetricia \quad [ \phantom{X} ] Odontología
\end{itemize}

\vspace{0.3cm}
\noindent\textbf{II. Preguntas de Apertura y Clima de Confianza (\textit{Rapport} y \textit{Napaykuy}):}
\begin{enumerate}
    \item ¿Podría relatarme qué motivo de salud le motivó a acudir hoy al Centro de Salud Belén y cómo se sentía anímicamente antes de ingresar al consultorio?\\
    \textit{En quechua chanka:} ¿Imataq qamta sasachasunki mamay/taytay, imaynam sonqoyki kachkarqa doctorman manaraq yaykuchkaspayki?
    \item ¿Es la primera vez que se atiende en este establecimiento o suele venir frecuentemente a este servicio?\\
    \textit{En quechua chanka:} ¿Qallariykutichu kay Centro de Saludpi atendichikunki, icha sapa kutichu hamunki?
\end{enumerate}

\vspace{0.3cm}
\noindent\textbf{III. Preguntas por Dimensiones Existenciales Fenomenológicas (van Manen y Sampieri):}
\begin{itemize}
    \item \textbf{Dimensión Existencial 1: Temporalidad (El tiempo vivido)}
    \begin{enumerate}[resume]
        \item ¿A qué hora tuvo que salir de su casa y a qué hora llegó a la fila para alcanzar un cupo? ¿Cómo vivió subjetivamente el transcurrir de las horas en la espera de la madrugada hasta culminar su atención?\\
        \textit{En quechua chanka:} ¿Ima oratataq wasikimanta lloqsimurqanki hinaspa ima oratataq chayamurqanki colaman cupota aypanaykipaq? ¿Imaynatam unay tiempota suyaspayki sientekurqanki lloqsinaykikama?
    \end{enumerate}

    \item \textbf{Dimensión Existencial 2: Espacialidad (El espacio vivido)}
    \begin{enumerate}[resume]
        \item ¿Cómo describe los espacios por donde transitó (la acera exterior a oscuras, la estrechez de los pasillos, las bancas de espera y el consultorio)? ¿Sintió que el ambiente de la consulta le brindó privacidad y tranquilidad?\\
        \textit{En quechua chanka:} ¿Imaynatam suyana sitiota qawarqanki (tutayaqpi calleta, pasadizokunata, bancakunata)? ¿Consultoriopi hawka kaytachu hinaspa secretotachu sientekurqanki rimanaykipaq?
    \end{enumerate}

    \item \textbf{Dimensión Existencial 3: Corporalidad (El cuerpo vivido)}
    \begin{enumerate}[resume]
        \item A lo largo de la madrugada y durante la consulta, ¿qué sensaciones experimentó físicamente su cuerpo (el frío penetrante, el dolor de su dolencia, el cansancio por estar de pie, o el alivio al ser atendido)?\\
        \textit{En quechua chanka:} ¿Suyasqayki horapi, imaynam cuerpoyki sientekurqa (chiriwan, nanaywan, sayasqa shaykuywan, icha doctorta rikuspa hawkayaywan)?
    \end{enumerate}

    \item \textbf{Dimensión Existencial 4: Relacionalidad (La relación humana vivida e intersubjetividad)}
    \begin{enumerate}[resume]
        \item ¿Cómo fue el encuentro humano con el médico, la enfermera o el personal de ventanilla? ¿Sintió calidez en la mirada, paciencia y escucha atenta? ¿Le hablaron en quechua chanka cuando usted lo requirió?\\
        \textit{En quechua chanka:} ¿Imaynam runa kayninkuwan chaskisurqanki doctor o enfermera? ¿Kuyakuywanchu qawasurqanki, pacienciawanchu uyarirqasunki? ¿Runasimipichu rimapayasurqanki mañakuptiki?
    \end{enumerate}

    \item \textbf{Dimensión 5: Síntesis de la Esencia y Significado Holístico del Cuidado Humano}
    \begin{enumerate}[resume]
        \item Al reflexionar sobre toda su experiencia, ¿qué significa para usted recibir una atención de auténtica calidad y dignidad, y qué propondría para humanizar el cuidado en el Centro de Salud Belén?\\
        \textit{En quechua chanka:} ¿Tukuy kaykunata yuyarispa, imataq qampaq allin digno atención kanman? ¿Imatataq mañakuwaq Centro de Saludpi aswan kuyakuywan atiendenankupaq?
    \end{enumerate}
\end{itemize}

\vspace{0.5cm}

% ------------------------------------------------------------------------------
% ANEXO 2: MATRIZ DE CONSISTENCIA CUALITATIVA FENOMENOLÓGICA
% ------------------------------------------------------------------------------
\phantomsection
\subsection*{Anexo 2: Matriz de Consistencia Fenomenológica Cualitativa}
\label{anx:2_matriz}
\addcontentsline{toc}{subsection}{Anexo 2: Matriz de Consistencia Cualitativa}

\noindent\textbf{Título:} PERCEPCIÓN DE LA CALIDAD DE ATENCIÓN EN USUARIOS DEL CENTRO DE SALUD BELÉN, AYACUCHO 2026.

\begin{table}[H]
\centering
\scriptsize
\caption{Matriz de consistencia cualitativa fenomenológica (EN 486 - UNSCH).}
\label{tab:matriz_consistencia}
\vspace{0.2cm}
\begin{tabularx}{\textwidth}{p{2.6cm}p{2.7cm}p{2.6cm}p{2.6cm}X}
\toprule
\textbf{Problema Fenomenológico} & \textbf{Objetivos Cualitativos} & \textbf{Preguntas Norteadoras} & \textbf{Dimensiones Existenciales} & \textbf{Metodología Cualitativa} \\
\midrule
\textbf{General:}\newline ¿Cuál es el significado, la estructura y la esencia de la experiencia vivida por los usuarios respecto al fenómeno del acceso y la calidad de atención en el Centro de Salud Belén, Ayacucho 2026? & 
\textbf{General:}\newline Explorar, comprender y describir el significado, estructura y esencia de la experiencia vivida por los usuarios respecto al acceso y calidad de atención en el Centro de Salud Belén, Ayacucho 2026. & 
1. Temporalidad (tiempo vivido en madrugadas a las 4:00 AM y esperas prolongadas).\newline 
2. Espacialidad (espacio vivido en veredas, pasillos y privacidad en consultorio).\newline 
3. Corporalidad (cuerpo vivido, frío de 5~°C a 8~°C, fatiga somática y dolor).\newline 
4. Relacionalidad (vínculo humano, empatía, acogida y quechua chanka).\newline 
5. Esencia compartida y variaciones individuales de la calidad asistencial. & 
\textbf{Categoría Central:}\newline Percepción del Fenómeno de Calidad de Atención y Acceso.\newline\newline 
\textbf{Existenciales del Lebenswelt:}\newline 
$\bullet$ D1: Temporalidad\newline 
$\bullet$ D2: Espacialidad\newline 
$\bullet$ D3: Corporalidad\newline 
$\bullet$ D4: Relacionalidad\newline 
$\bullet$ D5: Esencia y Divergencias & 
\textbf{Paradigma y Enfoque:}\newline Interpretativo-fenomenológico, naturalista (Husserl, van Manen, Sampieri).\newline\newline 
\textbf{Muestra e Informantes:}\newline Casos tipo no probabilístico. Criterio de saturación teórica (12 a 18 informantes clave).\newline\newline 
\textbf{Técnicas e Instrumentos:}\newline Entrevista fenomenológica en profundidad bilingüe, observación reflexiva y notas de campo.\newline\newline 
\textbf{Análisis y Rigor:}\newline Método de Colaizzi / van Manen con ATLAS.ti v9. Criterios de rigor de Guba y Lincoln. \\
\bottomrule
\end{tabularx}
\end{table}

\vspace{0.5cm}

% ------------------------------------------------------------------------------
% ANEXO 3: FORMATO DE JUICIO DE EXPERTOS CUALITATIVO
% ------------------------------------------------------------------------------
\phantomsection
\subsection*{Anexo 3: Formato de Validación Cualitativa por Juicio de Expertos}
\label{anx:3_expertos}
\addcontentsline{toc}{subsection}{Anexo 3: Formato de Juicio de Expertos}

\noindent\textbf{I. ASPECTOS GENERALES:}\\
1.1. Apellidos y nombres del informante (Experto): \hrulefill \\
1.2. Grado académico del experto: \hrulefill \\
1.3. Profesión del experto: \hrulefill \\
1.4. Institución donde labora: \hrulefill \\
1.5. Cargo que desempeña: \hrulefill \\
1.6. Denominación del instrumento: Guía de entrevista a profundidad fenomenológica semiestructurada bilingüe.\\
1.7. Autora del instrumento: Gresly Lucero Pariona Palomino.\\
1.8. Título de la tesis: Percepción de la calidad de atención en usuarios del Centro de Salud Belén, Ayacucho 2026.

\vspace{0.3cm}
\noindent\textbf{II. CRITERIOS DE RIGOR Y VALIDACIÓN CUALITATIVA (SAMPIERI, 2014):}

\begin{table}[H]
\centering
\small
\begin{tabularx}{\textwidth}{p{3.2cm}Xccp{3.5cm}}
\toprule
\textbf{Criterios} & \textbf{Indicadores de Evaluación Cualitativa} & \textbf{Sí} & \textbf{No} & \textbf{Observaciones} \\
\midrule
1. Claridad dialógica & Preguntas formuladas con lenguaje coloquial, abierto y comprensible para el usuario andino. & [ \phantom{X} ] & [ \phantom{X} ] & \\
\addlinespace
2. Pertinencia lingüística & Preguntas en quechua chanka gramaticalmente naturales, afectuosas y culturalmente pertinentes. & [ \phantom{X} ] & [ \phantom{X} ] & \\
\addlinespace
3. Consistencia fenomenológica & Articulación rigurosa con los cuatro existenciales del mundo de la vida (van Manen y Sampieri). & [ \phantom{X} ] & [ \phantom{X} ] & \\
\addlinespace
4. Coherencia vivencial & Correspondencia simétrica entre las preguntas y la captura de la esencia y las divergencias vivenciales. & [ \phantom{X} ] & [ \phantom{X} ] & \\
\addlinespace
5. Relevancia del cuidado & Explora estrictamente la interacción humana, los sentires corporales y la dignidad del cuidado. & [ \phantom{X} ] & [ \phantom{X} ] & \\
\addlinespace
6. Capacidad de saturación & Capacidad inductora de las preguntas para suscitar relatos densos y alcanzar la saturación de categorías. & [ \phantom{X} ] & [ \phantom{X} ] & \\
\bottomrule
\end{tabularx}
\caption{Ficha técnica de validación cualitativa por juicio de expertos.}
\label{tab:juicio_expertos}
\end{table}

\noindent\textbf{Observaciones generales:}\\
\hrulefill \\[0.2cm]
\hrulefill \\[0.3cm]
Ayacucho, \line(1,0){30} de \line(1,0){60} del 2026.\\[1.2cm]
\begin{center}
    \line(1,0){200}\\
    \textbf{FIRMA Y SELLO DEL EXPERTO}
\end{center}

\vspace{0.5cm}

% ------------------------------------------------------------------------------
% ANEXO 4: CONSENTIMIENTO INFORMADO
% ------------------------------------------------------------------------------
\phantomsection
\subsection*{Anexo 4: Consentimiento Informado para Participantes Informantes Clave}
\label{anx:4_consentimiento}
\addcontentsline{toc}{subsection}{Anexo 4: Modelo de Consentimiento Informado}

\noindent Yo, \line(1,0){220}, identificado con DNI N° \line(1,0){80}, domiciliado en \line(1,0){200}, distrito de Ayacucho, provincia de Huamanga.\\

\noindent He sido invitado(a) a participar de manera libre y voluntaria en la investigación titulada: \textbf{``PERCEPCIÓN DE LA CALIDAD DE ATENCIÓN EN USUARIOS DEL CENTRO DE SALUD BELÉN, AYACUCHO 2026''}, desarrollada por la bachiller Gresly Lucero Pariona Palomino, bajo la asesoría del Dr. Manglio Aguirre Andrade de la Universidad Nacional de San Cristóbal de Huamanga (UNSCH).\\

\noindent Declaro que se me han brindado las siguientes garantías de forma clara en mi idioma materno:
\begin{enumerate}
    \item Mi participación consistirá en compartir mis experiencias y sentires personales mediante una entrevista dialógica.
    \item La sesión será grabada únicamente en audio con fines de transcripción fidedigna.
    \item Mi identidad se mantendrá en estricta reserva y anonimato mediante el empleo de un código alfanumérico confidencial.
    \item Mi participación es 100\% voluntaria; puedo suspender la entrevista o retirarme del estudio en el momento que desee sin que ello perjudique en absoluto la atención médica que recibo en el establecimiento.
\end{enumerate}

\noindent En señal de plena conformidad, suscribo el presente documento y estampo mi huella dactilar.\\

\noindent Ayacucho, \line(1,0){30} de \line(1,0){80} del 2026.\\[1.2cm]
\begin{center}
\begin{minipage}{0.45\textwidth}
    \centering
    \line(1,0){150}\\
    \textbf{Firma del Participante Informante}\\
    DNI N°: \line(1,0){80}
\end{minipage}
\hfill
\begin{minipage}{0.45\textwidth}
    \centering
    \begin{tabular}{|c|}
        \hline
        \rule{0pt}{2.5cm}\parbox[b]{2.5cm}{\centering \scriptsize Huella Dactilar} \\
        \hline
    \end{tabular}
\end{minipage}
\end{center}

\vspace{0.5cm}

% ------------------------------------------------------------------------------
% ANEXO 5: CARTA DE ASESORÍA
% ------------------------------------------------------------------------------
\phantomsection
\subsection*{Anexo 5: Carta de Aceptación y Asesoría Formal del Proyecto de Tesis}
\label{anx:5_carta}
\addcontentsline{toc}{subsection}{Anexo 5: Carta de Asesoría Formal}

\noindent Ayacucho, 15 de enero del 2026.\\[0.3cm]
\textbf{Dr. Alejandro Yarlequé Mujica}\\
Decano de la Facultad de Ciencias de la Salud\\
Universidad Nacional de San Cristóbal de Huamanga\\[0.3cm]
\textbf{Asunto: Aceptación y Compromiso de Asesoría de Proyecto de Tesis Cualitativo}\\[0.3cm]
De mi mayor consideración:\\[0.2cm]
Tengo el agrado de dirigirme a su digno despacho para hacer de su conocimiento que he aceptado formalmente brindar la asesoría académica y metodológica al Proyecto de Tesis con enfoque cualitativo y diseño fenomenológico titulado:
\begin{center}
\textbf{``PERCEPCIÓN DE LA CALIDAD DE ATENCIÓN EN USUARIOS DEL CENTRO DE SALUD BELÉN, AYACUCHO 2026''}
\end{center}
presentado por la egresada \textbf{Gresly Lucero Pariona Palomino}, para optar el Título Profesional de Licenciada en Enfermería, en el marco de la cátedra de Proyecto de Investigación en Salud (EN 486).\\[0.2cm]
Dicha labor de asesoría se desarrollará garantizando el más riguroso estándar científico y bioético durante las fases de formulación, dictamen ético, inmersión en campo, análisis fenomenológico, redacción del informe final y sustentación pública.\\[0.2cm]
Sin otro particular, reitero a usted los sentimientos de mi especial estima personal.\\[0.3cm]
Atentamente,\\[1.2cm]
\begin{center}
    \line(1,0){220}\\
    \textbf{Dr. Manglio Aguirre Andrade}\\
    Docente de la Facultad de Ciencias de la Salud -- UNSCH\\
    ORCID: 0000-0001-8234-567X \quad | \quad DNI: \line(1,0){60}
\end{center}
